\begin{tikzpicture}
    \node[font=\small\bfseries] at (5,1.2) {Laboratorio: colisión frontal};
    \node[particula] (electronDoce) at (0,0) {$e^-$};
    \node[particula] (fotonDoce) at (10,0) {$\gamma$};
    \draw[movimiento] (electronDoce.east) -- (4,0) node[midway,above] {$+P_e$};
    \draw[movimiento] (fotonDoce.west) -- (7.8,0) node[midway,above] {$-E_\gamma/c$};
    \node[below=5mm] at (0,0) {$E_e=50\ \mathrm{MeV}$};
    \node[below=5mm] at (10,0) {$E_\gamma=25\ \mathrm{MeV}$};
    \draw[eje] (0,-1.3) -- (11,-1.3) node[right] {$x$};
    \draw[movimiento] (3,-2) -- (7,-2) node[midway,above] {$V_{\mathrm{CM}}\approx c/3$};
    \node at (5,-2.6) {El centro de momento se mueve hacia la derecha};
\end{tikzpicture}

{\small Se toma $+x$ en la dirección del electrón. Los momentos incidentes son opuestos, pero el del electrón tiene mayor magnitud.}